% tikz-style-cn-review.tex —— 中文文档「ASCII 框图 → TikZ」共享样式（毛概青/玫红）
% 用法：在正文 preamble \input 本文件（或把其中 \tikzset/\definecolor 复制进去）；
%       ASCII→TikZ 转换出的每张图用其中的 rvmod/rvbase/rvaccent/rvann/rvarr 等样式。
% 引擎：XeLaTeX / Tectonic（tcolorbox 会载入 tikz/pgf）。
\usepackage[most]{tcolorbox}
\usetikzlibrary{positioning,arrows.meta,fit,backgrounds,calc,shapes.geometric,shapes.arrows,chains}
\definecolor{rvteal}{HTML}{2F5C6E}
\definecolor{rvrose}{HTML}{C0504D}
\definecolor{rvtealbg}{HTML}{EAF1F4}
\definecolor{rvrosebg}{HTML}{FBEEEE}
\definecolor{rvgray}{HTML}{6B6B6B}
\tikzset{
  rvmod/.style   ={draw=rvteal, fill=rvtealbg, rounded corners=3pt, line width=0.6pt, align=center, inner sep=4pt, font=\small},
  rvbase/.style  ={draw=rvteal, fill=rvteal, text=white, rounded corners=4pt, line width=0.8pt, align=center, inner sep=5pt, font=\small},
  rvaccent/.style={draw=rvrose, fill=rvrosebg, rounded corners=3pt, line width=0.6pt, align=center, inner sep=4pt, font=\small},
  rvplain/.style ={draw=rvgray!55, fill=white, rounded corners=3pt, line width=0.5pt, align=center, inner sep=4pt, font=\small},
  rvann/.style   ={rvrose, font=\footnotesize\bfseries, align=left, inner sep=1pt},
  rvarr/.style   ={-{Stealth[length=2.4mm]}, rvteal, line width=0.8pt},
  rvline/.style  ={rvteal, line width=0.8pt},
}
% 图题：\rvcap{标题}
\newcommand{\rvcap}[1]{{\small\bfseries\color{rvteal} #1}\\[1.5mm]}
% 语义：rvmod=普通模块框 · rvbase=主干/容器标题(实心青白字) · rvaccent=强调(玫红框)
%       rvplain=中性灰框 · rvann=玫红"←说明"批注 · rvarr=青箭头 · rvline=青线
